\subsection*{\texttt{\_\_device\_\_} 函数：与 Croqtile 一并面向目标平台}

自定义归约、排序网络和特殊数学运算等计算，可以通过设备辅助函数实现。

\co@__co__@ 函数可以用 \co@call@ 关键字调用 \co@__device__@ 函数：

\begin{lstlisting}[language=croqtile,caption={device 函数：与 Croqtile 一并面向目标平台}]
#define N_TOKENS 1024
#define N_EXPERTS 8
#define N_THREAD 128
template <int K>
__device__ void warp_topk(f32* vals, s32* idxs);
template <typename T>
__device__ __forceinline__ T SHFL_XOR(T var, int lane_mask, int width) {
  return __shfl_xor_sync(uint32_t(-1), var, lane_mask, width);
}
__co__ void moe_topk(f32 [N_TOKENS, N_EXPERTS] scores,
                     s32 [N_TOKENS, K]& topk_ids,
                     f32 [N_TOKENS, K]& topk_scores,
                     int N_BLOCK) {
  parallel n by N_BLOCK : block {
    foreach m in [ |scores.span| / N_THREAD / N_BLOCK ] {
      shared_scores = dma.copy scores.chunkat(n#m, _) => shared;
      parallel gid by [N_THREAD / 32] : group
        parallel tid by 32 : thread {
          flat_scores = shared_scores.span_as([|shared_scores.span|]);
          score = flat_scores.at(gid # tid);
          local s32 [K] frag_idx{-1};
          local f32 [K] frag_val{-1.0f};
          frag_idx.at(0) = tid;
          frag_val.at(0) = score;
          call warp_topk<8>(frag_val, frag_idx);
          inthreads (tid == 0)
            foreach k in K {
              topk_ids.at(n#m, gid#k) = frag_idx.at(k);
              topk_scores.at(n#m, gid#k) = frag_val.at(k);
            }
        }
    }
  }
}
\end{lstlisting}

辅助函数接收指向元素存储的原始指针。因此，接口需要给出它所需的尺寸或布局约定，仅凭指针不会保留 Croqtile 形状。

\co@SHFL_XOR@ 辅助函数是交给目标编译器的普通设备代码。计算中如何使用它属于辅助函数的实现，Croqtile 不会把其函数体翻译成分块操作。

语句 \co@call group_topk<8>(...)@ 提供辅助函数的编译时配置和普通实参。Croqtile 将带形状数据实参转换为元素指针，由目标编译器验证生成的 C++ 调用。

区域 \co@inthreads (tid == 0)@ 选择一个线程完成其中工作，结束后汇合。它没有 \co@.async@ 修饰符，因此不会引入独立推进的生产者或消费者。

这种分工——Croqtile 管并行、分块和内存编排，设备函数管专门计算——在 MoE（专家混合）top-K、自定义归约和专用数学运算等生产核函数里很常见。
